\section{Intersection coordinates, endpoint jet, and restart}

On every strict horizon, the compatible intersection and the rectangle
relations arise independently from the same mixed jet. Theorems
\ref{thm:datum-intersection} and \ref{thm:datum-rectangles} carry those
two constructions to the whole terminal interval of the same forced history. We
first identify how the finite records represent \(\cI[u_0,\nu,f]\), then use
their adjacent rows to obtain a common endpoint and restart that history with
the same prescribed force.


\begin{definition}[Finite-coordinate representation of the intersection]
\label{def:compatible-intersection}
Let
\begin{equation}
 \mathscr X_\infty
 :=\varprojlim_{N,M}\mathscr X_{N,M}
 =\left\{(X_{N,M})_{N,M}:
 \begin{aligned}
 &X_{N,M}\in\mathscr X_{N,M},\\
 &\rho_{N',M'}^{N,M}X_{N,M}=X_{N',M'}
 \quad\text{whenever }N'\le N,\ M'\le M
 \end{aligned}
 \right\}.
 \label{eq:compatible-intersection-space}
\end{equation}
For a proposed finite maximal time, the canonical coordinate representation
of Theorem~\ref{thm:datum-intersection} identifies the whole-terminal
\(\cI[u_0,\nu,f]\) with
\begin{equation}
 \cI[u_0,\nu,f]
 \simeq\bigl(\pi_{N,M}\cI[u_0,\nu,f]\bigr)_{N,M\in\Nzero}
 =\bigl(\cR_{N,M}[u_0,\nu,f;T_*]\bigr)_{N,M\in\Nzero}.
 \label{eq:compatible-intersection}
\end{equation}
\end{definition}

\Needspace{14\baselineskip}
\begin{theorem}[Compatible finite-coordinate representation]
\label{thm:compatible-inverse-intersection}

Let \(\nu>0\), \(u_0\in\HsSigma(\R^3)\), prescribe
\(f\in C^\infty([0,\infty);\mathcal S(\R^3;\R^3))\), and let
\((T_*,u,p)\) be the maximal pressure-complete classical history generated
by \((u_0,\nu,f)\). If \(T_*<\infty\), then
\begin{equation}
 \cI[u_0,\nu,f]
 \simeq\bigl(\cR_{N,M}[u_0,\nu,f;T_*]\bigr)_{N,M\in\Nzero}
 \in\mathscr X_\infty
 \label{eq:ier-compatible-family}
\end{equation}
and all of its finite coordinates are the restrictions of the original
pressure-complete history.  In this representation, the coordinates
\((V_n,Q_n,F_n)_{n\in\Nzero}\) retain the memberships stated in
Theorem~\ref{thm:datum-intersection}:
\begin{align}
 V_n&\in\bigcap_{m\in\Nzero}L^2(0,T_*;H^m_\sigma(\R^3)),
 \label{eq:ier-inverse-velocity-coordinate}\\
 Q_n&\in\bigcap_{m\in\Nzero}L^1(0,T_*;H^m(\R^3)),
 \label{eq:ier-inverse-pressure-coordinate}\\
 F_n&\in\bigcap_{m\in\Nzero}C([0,T_*];H^m(\R^3)).
 \label{eq:ier-inverse-force-coordinate}
\end{align}
Every finite restriction is the corresponding rectangle
\(\cR_{N,M}[u_0,\nu,f;T_*]\).  The velocity membership inherited from that
theorem is
\begin{equation}
 u\in\bigcap_{r,m\in\Nzero}H^r(0,T_*;H^m(\R^3)),
 \label{eq:mixed-sobolev-intersection}
\end{equation}
and, for every \(r,m\in\Nzero\),
\begin{equation}
 \sum_{n=0}^{r}\norm{\partial_t^nu}_{L^2(0,T_*;H^m)}^2
 \le\mathfrak R_{r,m}^*[u_0,\nu,f,T_*].
 \label{eq:mixed-sobolev-finite-bound}
\end{equation}
\end{theorem}

\begin{proof}
The coordinate memberships
\eqref{eq:ier-inverse-velocity-coordinate},
\eqref{eq:ier-inverse-pressure-coordinate}, and
\eqref{eq:mixed-sobolev-intersection} are inherited from
Theorem~\ref{thm:datum-intersection}.
The membership \eqref{eq:ier-inverse-force-coordinate} follows directly from
the prescribed class
\(f\in C^\infty([0,\infty);\mathcal S(\R^3;\R^3))\).
For each finite \(N,M\), apply the projection \(\pi_{N,M}\).
Its entries are those of the rectangle selected directly from the mixed
jet.
Theorem~\ref{thm:datum-rectangles} and
Proposition~\ref{prop:rectangle-pressure-completion} give its finite norm
bounds.  Proposition~\ref{prop:rectangle-compatibility} gives
\[
 \rho_{N',M'}^{N,M}\pi_{N,M}\cI[u_0,\nu,f]
 =\pi_{N',M'}\cI[u_0,\nu,f]
 \qquad(N'\le N,\ M'\le M).
\]
Hence these projections form an element of \(\mathscr X_\infty\).
Every \(V_n,Q_n,F_n\) appears in a sufficiently deep projection, and every
spatial derivative is fixed by its distributional velocity, pressure, or force
row.  The full family consequently identifies precisely
\(\cI[u_0,\nu,f]\); this is a coordinate representation of that object, not
another construction of its existence.

Fix finite \(r,m\), put \(I=(0,T_*)\), and choose \(N\ge r\),
\(M\ge m\). The upper rows of that finite record belong to
\(L^2(I;H^{M+1})\hookrightarrow L^2(I;H^m)\) for \(0\le n\le r\).
Proposition~\ref{prop:formal-mixed-recurrences} identifies them as the
distributional derivatives \(V_n=\partial_t^nu\). Thus the finite record
contains every term in the derivative-sum norm of \(H^r(I;H^m)\).
Taking \(N=r\), \(M=m\) gives
\[
 \begin{aligned}
 \norm u_{H^r(I;H^m)}^2
 &=\sum_{n=0}^r\norm{V_n}_{L^2(I;H^m)}^2\\
 &\le\sum_{n=0}^r\norm{V_n}_{L^2(I;H^{m+1})}^2
 \le\mathfrak R_{r,m}^*[u_0,\nu,f,T_*].
 \end{aligned}
\]
This proves \eqref{eq:mixed-sobolev-finite-bound} and shows how each
mixed Sobolev norm is recovered from a finite record.  Its whole-terminal
membership is supplied by Theorem~\ref{thm:datum-intersection}; its finite
complete-equation estimate is supplied independently by
Theorem~\ref{thm:datum-rectangles}.
\end{proof}

\begin{proposition}[Complete spatial row and common endpoint]
\label{prop:ier-complete-spatial-endpoint}
Let the prescribed datum and finite maximal history be as in
Theorem~\ref{thm:compatible-inverse-intersection}.  Then, for every
\(m\in\Nzero\),
\begin{equation}
 u\in W^{1,1}(0,T_*;H^m_\sigma(\R^3))
 \hookrightarrow C([0,T_*];H^m_\sigma(\R^3)).
 \label{eq:ier-spatial-row-W11}
\end{equation}
The endpoint values supplied at the different Sobolev orders are the images
of one common state
\begin{equation}
 u_*:=u(T_*)\in\bigcap_{m\in\Nzero}H^m_\sigma(\R^3)
 =H^\infty_\sigma(\R^3).
 \label{eq:ier-common-spatial-endpoint}
\end{equation}
\end{proposition}

\begin{proof}
Fix \(m\).  The zeroth temporal coordinate of
\eqref{eq:mixed-sobolev-intersection} gives
\[
 u\in L^2(0,T_*;H^{m+2}).
\]
The pressure-completed nonlinearity satisfies
\[
 \norm{\Pp\nabla\!\cdot(u\otimes u)}_{H^m}
 \le C_m\norm u_{H^{m+2}}^2
\]
by Lemmas~\ref{lem:sobolev-product} and \ref{lem:riesz-leray}.  It belongs
to \(L^1(0,T_*;H^m)\).  This spatial row and \(T_*<\infty\) give
 \(\Delta u\in L^1(0,T_*;H^m)\), and the projected equation
\[
 u_t=\nu\Delta u-\Pp\nabla\!\cdot(u\otimes u)+\Pp f
\]
places \(u_t\) in \(L^1(0,T_*;H^m)\), because the prescribed source is
integrable there in every finite Sobolev order.  Also
\(u\in L^1(0,T_*;H^m)\), so the Banach-valued fundamental theorem gives
\eqref{eq:ier-spatial-row-W11}.

If \(k\ge m\), the \(H^k\)- and \(H^m\)-continuous representatives are
the original preterminal function after the embedding \(H^k\hookrightarrow
H^m\).  Uniqueness of continuous representatives makes their endpoint
values agree under that embedding.  These compatible values define the
single state in \eqref{eq:ier-common-spatial-endpoint}.  Continuity of
divergence from \(H^{m+1}\) to \(H^m\) gives \(\nabla\cdot u_*=0\).
\end{proof}

\begin{corollary}[Whole-terminal pressure-complete mixed jet]
\label{cor:ier-whole-terminal-mixed-jet}
Let the prescribed datum and finite maximal history be as in
Theorem~\ref{thm:compatible-inverse-intersection}.  Then, for every
\(r,m,n\in\Nzero\) and \(\alpha\in\Nzero^3\),
\begin{align}
 u&\in H^r(0,T_*;H^m_\sigma(\R^3)),
 \label{eq:ier-whole-terminal-bochner}\\
 J_{n,\alpha}:=\partial_x^\alpha V_n
 &\in L^2(0,T_*;H^{m+1}(\R^3)),
 \label{eq:ier-whole-terminal-mixed-upper}\\
 J_{n+1,\alpha}
 &\in L^2(0,T_*;H^{m-1}(\R^3)),
 \label{eq:ier-whole-terminal-mixed-lower}\\
 \Pi_{n,\alpha}:=\partial_x^\alpha Q_n
 &\in L^1(0,T_*;H^m(\R^3)).
 \label{eq:ier-whole-terminal-mixed-pressure}
\end{align}
The following finite bounds hold:
\begin{align}
 \sum_{j=0}^{r}\norm{V_j}_{L^2(0,T_*;H^m)}^2
 &\le\mathfrak R_{r,m}^*[u_0,\nu,f,T_*],
 \label{eq:ier-finite-bochner-bound}\\
 \norm{J_{n,\alpha}}_{L^2(0,T_*;H^{m+1})}^2
 &\quad+\norm{J_{n+1,\alpha}}_{L^2(0,T_*;H^{m-1})}^2
 \notag\\
 &\le\mathfrak R_{n,m+|\alpha|}^*[u_0,\nu,f,T_*],
 \label{eq:ier-finite-mixed-bound}\\
 \norm{\Pi_{n,\alpha}}_{L^1(0,T_*;H^m)}
 &\le C_{m+|\alpha|}^{\rm p}(2^{n+1}-1)
 \mathfrak R_{n,m+|\alpha|}^*[u_0,\nu,f,T_*]
 +\mathfrak F^{\rm p}_{n,\alpha,m}(0,T_*).
 \label{eq:ier-finite-mixed-pressure-bound}
\end{align}
Every mixed row satisfies the recurrences in
Proposition~\ref{prop:formal-mixed-recurrences} on \((0,T_*)\).
\end{corollary}

\begin{proof}
Equations \eqref{eq:ier-whole-terminal-bochner} and
\eqref{eq:ier-finite-bochner-bound} follow from
Theorem~\ref{thm:compatible-inverse-intersection}.  The Fourier multiplier
bound
\[
 \norm{\partial_x^\alpha h}_{H^s}
 \le\norm h_{H^{s+|\alpha|}}
\]
and the rectangle \((n,m+|\alpha|)\) give
\eqref{eq:ier-whole-terminal-mixed-upper}--%
\eqref{eq:ier-finite-mixed-bound}.  Apply this multiplier bound to the
pressure completion in Proposition~\ref{prop:rectangle-pressure-completion}
to obtain \eqref{eq:ier-whole-terminal-mixed-pressure} and
\eqref{eq:ier-finite-mixed-pressure-bound}.  Proposition
\ref{prop:formal-mixed-recurrences} supplies the stated identities.
\end{proof}

\Needspace{24\baselineskip}
\begin{theorem}[Compatible whole-terminal Sobolev traces]
\label{thm:ier-compatible-endpoint-traces}
Let the prescribed datum and finite maximal history be as in
Theorem~\ref{thm:compatible-inverse-intersection}.
For \(N,m\in\Nzero\), define
\begin{equation}
 \mathfrak E_{N,m}
 :=(1+T_*^{-1})^{1/2}
 \bigl(\mathfrak R_{N,m}^*[u_0,\nu,f,T_*]\bigr)^{1/2}.
 \label{eq:ier-endpoint-majorant}
\end{equation}
\begin{samepage}
If \(0\le n\le N\) and \(\alpha\in\Nzero^3\), then
\begin{align}
 V_n&\in C([0,T_*];H^m_\sigma(\R^3)),
 &\sup_{0\le t\le T_*}\norm{V_n(t)}_{H^m}
 &\le\mathfrak E_{N,m},
 \label{eq:ier-row-trace}\\
 J_{n,\alpha}&\in C([0,T_*];H^m(\R^3)),
 &\sup_{0\le t\le T_*}\norm{J_{n,\alpha}(t)}_{H^m}
 &\le\mathfrak E_{N,m+|\alpha|}.
 \label{eq:ier-mixed-row-trace}
\end{align}
There are unique values
\begin{align}
 V_n^*&:=\lim_{t\uparrow T_*}V_n(t)
 \in H^\infty(\R^3),
 \label{eq:ier-velocity-endpoint}\\
 J_{n,\alpha}^*&:=\lim_{t\uparrow T_*}J_{n,\alpha}(t)
 =\partial_x^\alpha V_n^*
 \in H^\infty(\R^3),
 \label{eq:ier-mixed-velocity-endpoint}
\end{align}
\par
\end{samepage}
such that, for \(k\ge m\), the image of the \(H^k\) trace under
\(H^k\hookrightarrow H^m\) is the \(H^m\) trace.  Quantitatively,
\begin{align}
 \norm{V_n^*}_{H^m}
 &\le\mathfrak E_{n,m},
 \label{eq:ier-endpoint-velocity-size}\\
 \norm{V_n(t)-V_n^*}_{H^m}
 &\le(T_*-t)\mathfrak E_{n+1,m},
 \qquad0\le t\le T_*,
 \label{eq:ier-endpoint-velocity-modulus}\\
 \norm{J_{n,\alpha}^*}_{H^m}
 &\le\mathfrak E_{n,m+|\alpha|},
 \label{eq:ier-endpoint-mixed-size}\\
 \norm{J_{n,\alpha}(t)-J_{n,\alpha}^*}_{H^m}
 &\le(T_*-t)\mathfrak E_{n+1,m+|\alpha|}.
 \label{eq:ier-endpoint-mixed-modulus}
\end{align}
 In particular, the zeroth trace agrees with the common endpoint in
 Proposition~\ref{prop:ier-complete-spatial-endpoint}, and
\begin{equation}
 u_*=V_0^*\in H^\infty_\sigma(\R^3),
 \qquad
 u\in\bigcap_{a,m\in\Nzero}C^a([0,T_*];H^m_\sigma(\R^3)),
 \qquad
 \partial_t^au(T_*)=V_a^*.
 \label{eq:ier-complete-left-velocity-jet}
\end{equation}
\end{theorem}

\begin{proof}
For \(0\le n\le N\), project the intersection
\(\cI[u_0,\nu,f]\) supplied by Theorem~\ref{thm:datum-intersection} onto the two
coordinates
\[
 V_n\in L^2(0,T_*;H^{m+1}),
 \qquad
 \partial_tV_n=V_{n+1}\in L^2(0,T_*;H^{m-1}).
\]
The complete mixed intersection supplies both finite norms.
With \(A=(1-\Delta)^{1/2}\), the adjacent-row identity is
\[
 \frac{\mathrm d}{\mathrm dt}\norm{V_n(t)}_{H^m}^2
 =2\operatorname{Re}\ip{A^{m-1}V_{n+1}(t)}{A^{m+1}V_n(t)}_{L^2}.
\]
Integrating from the initial time and applying Cauchy--Schwarz gives
\begin{equation}
 \sup_{0\le t\le T_*}\norm{V_n(t)}_{H^m}^2
 \le\norm{V_n(0)}_{H^m}^2
 +2\norm{V_{n+1}}_{L^2(I;H^{m-1})}
    \norm{V_n}_{L^2(I;H^{m+1})},
 \qquad I=(0,T_*).
 \label{eq:ier-initial-row-trace-bound}
\end{equation}
The initial norm is finite by Lemma~\ref{lem:formal-initial-jet}.
Lemma~\ref{lem:adjacent-trace} supplies the continuous representative
on \([0,T_*]\), justifies the identity by Fourier truncation, and gives the
first membership in
\eqref{eq:ier-row-trace} and
\begin{align*}
 \sup_{0\le t\le T_*}\norm{V_n(t)}_{H^m}^2
 &\le T_*^{-1}\norm{V_n}_{L^2H^m}^2
 +2\norm{V_n}_{L^2H^{m+1}}\norm{V_{n+1}}_{L^2H^{m-1}}\\
 &\le(1+T_*^{-1})\mathfrak R_{N,m}^*[u_0,\nu,f,T_*].
\end{align*}
This proves the bound in \eqref{eq:ier-row-trace}.  Apply this lemma to
\(J_{n,\alpha}\), using \eqref{eq:ier-finite-mixed-bound}, to prove
\eqref{eq:ier-mixed-row-trace}.

For \(k\ge m\), the \(H^k\)-continuous and \(H^m\)-continuous
representatives equal the original preterminal distribution.  After applying
\(H^k\hookrightarrow H^m\), uniqueness of continuous representatives makes
them equal on \([0,T_*]\).  Their endpoint values are compatible and define
\eqref{eq:ier-velocity-endpoint}--%
\eqref{eq:ier-mixed-velocity-endpoint}.  Spatial differentiation is
continuous between the corresponding finite Sobolev spaces, so
\(J_{n,\alpha}^*=\partial_x^\alpha V_n^*\).

Corollary~\ref{cor:ier-whole-terminal-mixed-jet} gives
\(V_n,V_{n+1}\in L^1(0,T_*;H^m)\) and
\(\partial_tV_n=V_{n+1}\) in distributions.  The Banach-valued fundamental
identity gives
\begin{equation}
 V_n(t)-V_n(s)=\int_s^tV_{n+1}(\tau)\dd\tau
 \quad\text{in }H^m,
 \qquad0\le s\le t\le T_*.
 \label{eq:ier-banach-fundamental-identity}
\end{equation}
The continuous representative of \(V_{n+1}\) is bounded by
\(\mathfrak E_{n+1,m}\).  Set \(t=T_*\) in
\eqref{eq:ier-banach-fundamental-identity} to prove
\eqref{eq:ier-endpoint-velocity-modulus}; spatial differentiation gives
\eqref{eq:ier-endpoint-mixed-modulus}.  The remaining size estimates follow
from \eqref{eq:ier-row-trace}--\eqref{eq:ier-mixed-row-trace}.

The divergence map \(H^{m+1}\to H^m\) is continuous, so
\(\nabla\cdot u_*=0\).  For each \(a,m\), the curves
\(V_0,\ldots,V_{a+1}\) are continuous in \(H^m\), and
\eqref{eq:ier-banach-fundamental-identity} identifies the derivative of
\(V_j\) with \(V_{j+1}\).  Induction in \(j\) proves
\eqref{eq:ier-complete-left-velocity-jet}.
\end{proof}

\begin{theorem}[Pressure-complete endpoint jet]
\label{thm:ier-pressure-complete-endpoint}
\begin{samepage}
Let the prescribed datum and finite maximal history be as in
Theorem~\ref{thm:compatible-inverse-intersection}.
For \(n\in\Nzero\), define
\[
 \Phi_n(t):=(-\Delta)^{-1}\nabla\cdot F_n(t).
\]
The prescribed force class gives
\(\Phi_n\in C^\infty([0,T_*+1];H^s)\) for every finite \(n,s\).  Define
\begin{equation}
 Q_n^*:=R_iR_j\sum_{a=0}^{n}\binom na
 V_{a,i}^*V_{n-a,j}^*-\Phi_n(T_*).
 \label{eq:ier-terminal-pressure-row}
\end{equation}
Then \(Q_n^*\in H^\infty(\R^3)\), and, for every \(s\in\Nzero\),
\begin{align}
 \norm{Q_n^*}_{H^s}
 &\le9C_s^{\rm alg}2^n\mathfrak E_{n,s+1}^2
 +\norm{\Phi_n(T_*)}_{H^s},
 \label{eq:ier-terminal-pressure-size}\\
 \norm{Q_n(t)-Q_n^*}_{H^s}
 &\le9C_s^{\rm alg}2^{n+1}(T_*-t)\mathfrak E_{n+1,s+1}^2\\
 &\quad +(T_*-t)
 \sup_{0\le r\le T_*}\norm{\Phi_{n+1}(r)}_{H^s}.
 \label{eq:ier-terminal-pressure-modulus}
\end{align}
\end{samepage}
For \(\alpha\in\Nzero^3\),
\begin{equation}
 \Pi_{n,\alpha}^*:=\partial_x^\alpha Q_n^*
 =\lim_{t\uparrow T_*}\Pi_{n,\alpha}(t)
 \quad\text{in }H^m(\R^3)\text{ for every }m\in\Nzero.
 \label{eq:ier-terminal-mixed-pressure-row}
\end{equation}
More precisely, for every \(m\in\Nzero\),
\begin{align}
 \norm{\Pi_{n,\alpha}^*}_{H^m}
 &\le9C_{m+|\alpha|}^{\rm alg}2^n
 \mathfrak E_{n,m+|\alpha|+1}^2
 +\norm{\partial_x^\alpha\Phi_n(T_*)}_{H^m},
 \label{eq:ier-terminal-mixed-pressure-size}\\
 \norm{\Pi_{n,\alpha}(t)-\Pi_{n,\alpha}^*}_{H^m}
 &\le9C_{m+|\alpha|}^{\rm alg}2^{n+1}(T_*-t)
 \mathfrak E_{n+1,m+|\alpha|+1}^2\\
 &\quad +(T_*-t)\sup_{0\le r\le T_*}
 \norm{\partial_x^\alpha\Phi_{n+1}(r)}_{H^m}.
 \label{eq:ier-terminal-mixed-pressure-modulus}
\end{align}
The endpoint rows satisfy, for every \(n,q\in\Nzero\),
\begin{equation}
 V_{n+1}^*
 =\nu\Delta V_n^*
 -\sum_{a=0}^{n}\binom na(V_a^*\cdot\nabla)V_{n-a}^*
 -\nabla Q_n^*+F_n(T_*)
 \quad\text{in }H^q(\R^3).
 \label{eq:ier-terminal-velocity-row}
\end{equation}
Among families rooted at \(u_*\), carrying the prescribed values
\((F_n(T_*))_{n\in\Nzero}\), and satisfying
\eqref{eq:ier-terminal-pressure-row} and
\eqref{eq:ier-terminal-velocity-row} in every finite Sobolev space,
\((V_n^*,Q_n^*)_{n\in\Nzero}\) is the unique pressure-normalized endpoint
jet.  Moreover,
\begin{equation}
 p\in\bigcap_{a,m\in\Nzero}C^a([0,T_*];H^m(\R^3)),
 \qquad
 \partial_t^ap(T_*)=Q_a^*.
 \label{eq:ier-complete-left-pressure-jet}
\end{equation}
\end{theorem}

\begin{proof}
Lemmas~\ref{lem:sobolev-product} and \ref{lem:riesz-leray}, followed by
\eqref{eq:ier-row-trace}, give
\begin{align*}
 \norm{Q_n^*}_{H^s}
 &\le9C_s^{\rm alg}\sum_{a=0}^{n}\binom na
 \norm{V_a^*}_{H^{s+1}}\norm{V_{n-a}^*}_{H^{s+1}}
 +\norm{\Phi_n(T_*)}_{H^s}\\
 &\le9C_s^{\rm alg}2^n\mathfrak E_{n,s+1}^2
 +\norm{\Phi_n(T_*)}_{H^s}.
\end{align*}
This is \eqref{eq:ier-terminal-pressure-size}.
For each product in \(Q_n(t)-Q_n^*\), insert
\(V_a^*V_{n-a}(t)\), use
\eqref{eq:ier-endpoint-velocity-modulus}, and bound every retained factor by
\(\mathfrak E_{n+1,s+1}\).  The two product differences and
\(\sum_{a=0}^{n}\binom na=2^n\) prove
the product part of \eqref{eq:ier-terminal-pressure-modulus}.  Since
\(\partial_t\Phi_n=\Phi_{n+1}\),
\[
 \norm{\Phi_n(t)-\Phi_n(T_*)}_{H^s}
 \le (T_*-t)\sup_{0\le r\le T_*}\norm{\Phi_{n+1}(r)}_{H^s},
\]
which supplies the remaining term.  The Fourier multiplier inequality
\[
 \norm{\partial_x^\alpha h}_{H^m}
 \le\norm h_{H^{m+|\alpha|}}
\]
applied to the nonlinear Riesz-product component, together with the exact
derivative \(\partial_x^\alpha\Phi_n\) of the force potential, proves
\eqref{eq:ier-terminal-mixed-pressure-size} and
\eqref{eq:ier-terminal-mixed-pressure-modulus}.  The latter tends to zero
as \(t\uparrow T_*\), which proves
\eqref{eq:ier-terminal-mixed-pressure-row} in every \(H^m\).

Fix \(n,q\).  The trace of \(V_n\) in \(H^{q+2}\) carries the diffusion
term in Proposition~\ref{prop:formal-mixed-recurrences} to \(H^q\) at
\(T_*\).  The product estimate at order \(q+1\), using velocity convergence
in \(H^{q+2}\), carries every transport term to \(H^q\).  Equation
\eqref{eq:ier-terminal-pressure-modulus} at order \(q+1\) carries the pressure
gradient to \(H^q\).  The left side converges to \(V_{n+1}^*\) in \(H^q\),
and \(F_n(t)\to F_n(T_*)\) in \(H^q\), which proves
\eqref{eq:ier-terminal-velocity-row}.

The value \(V_0^*=u_*\) fixes \(Q_0^*\), and the velocity recurrence fixes
\(V_1^*\).  Inductively, \(V_0^*,\ldots,V_n^*\) fix \(Q_n^*\), which fixes
\(V_{n+1}^*\).  This proves uniqueness within the displayed recurrence
class.  Extend each \(Q_n\) to \(T_*\) by
\(Q_n(T_*):=Q_n^*\).  On \((0,T_*)\),
\(Q_n=\partial_t^np\), so \(\partial_tQ_n=Q_{n+1}\) in distributions.
Both rows are continuous in every \(H^m\) by
\eqref{eq:ier-terminal-pressure-modulus}.  Hence, for
\(0\le s\le t\le T_*\),
\begin{equation}
 Q_n(t)-Q_n(s)=\int_s^tQ_{n+1}(\tau)\dd\tau
 \quad\text{in }H^m(\R^3).
 \label{eq:ier-pressure-banach-identity}
\end{equation}
Indeed, the identity holds for \(t<T_*\), and the endpoint bounds permit
passage to \(t=T_*\).  Thus
\(Q_n\in C^1([0,T_*];H^m)\) with derivative \(Q_{n+1}\).  Induction gives
\eqref{eq:ier-complete-left-pressure-jet}.
\end{proof}

\begin{lemma}[Late-slice overlap uniqueness]
\label{lem:ier-late-slice-overlap}
Let the prescribed datum and finite maximal history be as in
Theorem~\ref{thm:compatible-inverse-intersection}.
Define
\[
 \Gamma_*:=\sup_{0\le s\le T_*}
 \norm{\Pp f}_{L^2(s,s+1;H^2)}<\infty.
\]
Define
\begin{equation}
 \delta_0
 :=\min\{1,c_\nu(1+\mathfrak E_{0,3}+\Gamma_*)^{-2}\}>0,
 \label{eq:ier-uniform-restart-time}
\end{equation}
where \(c_\nu\) is the constant in Theorem~\ref{thm:local-restart}.
For every
\(t_0\in(\max\{0,T_*-\delta_0/2\},T_*)\), let
\[
 (w^{t_0},q^{t_0})
 \in\mathscr M_{\nu,f}^3(t_0,t_0+\delta_0;u(t_0)),
 \qquad
 (w^*,q^*)
 \in\mathscr M_{\nu,f}^3(T_*,T_*+\delta_0;u_*)
\]
be the unique pairs supplied by Theorem~\ref{thm:local-restart}.  Then
\begin{align}
 (w^{t_0},q^{t_0})&=(u,p)
 &&\text{on }[t_0,T_*),
 \label{eq:ier-left-overlap}\\
 (w^{t_0},q^{t_0})&=(w^*,q^*)
 &&\text{on }[T_*,t_0+\delta_0].
 \label{eq:ier-right-overlap}
\end{align}
\end{lemma}

\begin{proof}
Equation \eqref{eq:ier-row-trace} gives
\[
 \norm{u(t_0)}_{H^3}\le\mathfrak E_{0,3},
 \qquad
 \norm{u_*}_{H^3}\le\mathfrak E_{0,3}.
\]
Hence
\[
 \delta_0\le\delta_{\nu,f}(t_0,u(t_0)),
 \qquad
 \delta_0\le\delta_{\nu,f}(T_*,u_*),
\]
with \(\delta_{\nu,f}\) as in Theorem~\ref{thm:local-restart}.
Theorem~\ref{thm:local-restart} gives both pairs on the stated intervals.
For each \(t_0<b<T_*\), Lemma~\ref{lem:shifted-mild-restriction} places
\((u,p)|_{[t_0,b]}\) and
\((w^{t_0},q^{t_0})|_{[t_0,b]}\) in
\(\mathscr M_{\nu,f}^3(t_0,b;u(t_0))\).  Uniqueness in that class gives
\eqref{eq:ier-left-overlap}.  Continuity in \(H^3\) gives
\[
 w^{t_0}(T_*)=\lim_{t\uparrow T_*}u(t)=u_*.
\]
Lemma~\ref{lem:shifted-mild-restriction} places the restrictions of
\((w^{t_0},q^{t_0})\) and \((w^*,q^*)\) to
\([T_*,t_0+\delta_0]\) in
\(\mathscr M_{\nu,f}^3(T_*,t_0+\delta_0;u_*)\).  Uniqueness in that class gives
\eqref{eq:ier-right-overlap}.  The right overlap has positive length because
\(t_0+\delta_0>T_*+\delta_0/2\).
\end{proof}

\begin{theorem}[Same-forced-history endpoint restart]
\label{thm:ier-same-history-restart}
Let the prescribed datum and finite maximal history be as in
Theorem~\ref{thm:compatible-inverse-intersection}.
There is a pressure-normalized forced classical solution, with the same
prescribed force \(f(t)\),
\((\widetilde u,\widetilde p)\) on
\([0,T_*+\delta_0/2]\times\R^3\) such that
\begin{align}
 (\widetilde u,\widetilde p)|_{[0,T_*)\times\R^3}
 &=(u,p),
 \label{eq:ier-extension-identity}\\
 \widetilde u,\widetilde p
 &\in C^\infty([0,T_*+\delta_0/2]\times\R^3),
 \label{eq:ier-extension-smoothness}\\
 \partial_t^n\widetilde u(T_*)&=V_n^*,
 \qquad
 \partial_t^n\widetilde p(T_*)=Q_n^*
 \quad(n\in\Nzero).
 \label{eq:ier-extension-jet}
\end{align}
If \(0<\eta\le\delta_0/2\) and
\((\widehat u,\widehat p)\) is another pressure-normalized classical
extension satisfying
\begin{equation}
 (\widehat u,\widehat p)|_{[0,T_*)}=(u,p),
 \qquad
 (\widehat u,\widehat p)|_{[T_*,T_*+\eta]}
 \in\mathscr M_{\nu,f}^3(T_*,T_*+\eta;u_*),
 \label{eq:ier-same-history-uniqueness-class}
\end{equation}
then
\begin{equation}
 (\widehat u,\widehat p)=(\widetilde u,\widetilde p)
 \quad\text{on }[0,T_*+\eta].
 \label{eq:ier-same-history-uniqueness}
\end{equation}
\end{theorem}

\begin{proof}
Let \((w^*,q^*)\) be the endpoint solution from
Lemma~\ref{lem:ier-late-slice-overlap}, and set
\[
 W_n^+:=\partial_t^nw^*(T_*),
 \qquad
 P_n^+:=\partial_t^nq^*(T_*).
\]
\begin{samepage}
The restart datum is exactly
\[
 w^*(T_*)=u_*=u(T_*).
\]
The equation continues with the same absolute-time source \(f(t)\).  If the
right-hand interval is parametrized by \(\tau=t-T_*\), that same source is
written \(f(T_*+\tau)\); this is only the shifted time coordinate.
Theorem~\ref{thm:local-restart} makes \((w^*,q^*)\) smooth and
pressure-normalized.  Differentiating its equations exactly as in the proof of
Proposition~\ref{prop:formal-mixed-recurrences} gives
\begin{align*}
 P_n^+&=R_iR_j\sum_{a=0}^{n}\binom naW_{a,i}^+W_{n-a,j}^+
 -(-\Delta)^{-1}\nabla\cdot F_n(T_*),\\
 W_{n+1}^+&=\nu\Delta W_n^+
 -\sum_{a=0}^{n}\binom na(W_a^+\cdot\nabla)W_{n-a}^+
 -\nabla P_n^++F_n(T_*).
\end{align*}
Since \(W_0^+=u_*=V_0^*\), induction with
\eqref{eq:ier-terminal-pressure-row} and
\eqref{eq:ier-terminal-velocity-row} gives
\begin{equation}
 W_n^+=V_n^*,
 \qquad P_n^+=Q_n^*,
 \qquad n\in\Nzero.
 \label{eq:ier-left-right-jet-match}
\end{equation}
\par
\end{samepage}

Define
\[
 (\widetilde u,\widetilde p)(t)
 :=\begin{cases}
 (u,p)(t),&0\le t<T_*,\\
 (w^*,q^*)(t),&T_*\le t\le T_*+\delta_0/2.
 \end{cases}
\]
Equations \eqref{eq:ier-complete-left-velocity-jet},
\eqref{eq:ier-complete-left-pressure-jet}, and
\eqref{eq:ier-left-right-jet-match} match every one-sided time derivative in
every finite Sobolev space.  Sobolev embedding gives
\eqref{eq:ier-extension-smoothness}, and the pressure-complete equation holds
at \(T_*\) by \eqref{eq:ier-terminal-velocity-row}.  Lemma
\ref{lem:ier-late-slice-overlap} identifies the joined pair near \(T_*\) with
the single local history \((w^{t_0},q^{t_0})\), so the gluing is one local
forced classical history across the endpoint; the piecewise definition gives
\eqref{eq:ier-extension-identity}.  Any pair in
\eqref{eq:ier-same-history-uniqueness-class} has that datum \(u_*\) on
the right of \(T_*\) and belongs to the shifted mild class containing
\((w^*,q^*)\).  The at-most-one assertion in
Theorem~\ref{thm:local-restart} proves
\eqref{eq:ier-same-history-uniqueness}.
\end{proof}

The datum-generated whole-terminal intersection of
Theorem~\ref{thm:datum-intersection}, projected through the independently
derived complete-equation rectangle identities by
Theorem~\ref{thm:datum-rectangles}, gives compatible velocity traces through \(T_*\) in
Theorem~\ref{thm:ier-compatible-endpoint-traces}.  The pressure
recurrence fixes their normalized endpoint jet in
Theorem~\ref{thm:ier-pressure-complete-endpoint}.  With those endpoint
values, Theorem~\ref{thm:ier-same-history-restart} extends the original
history beyond \(T_*\).  A proposed finite maximal time would admit a
classical continuation.  Corollary~\ref{cor:ier-maximal-time-contradiction}
applies this contradiction, and Theorem~\ref{thm:ier-global-smoothness}
collects the global solution, pressure normalization, and energy identity.

\begin{corollary}[Maximal-time contradiction for prescribed smooth data]
\label{cor:ier-maximal-time-contradiction}
Let \(\nu>0\), \(u_0\in\HsSigma(\R^3)\), prescribe
\(f\in C^\infty([0,\infty);\mathcal S(\R^3;\R^3))\), and let
\((T_*,u,p)\) be the maximal pressure-complete classical history generated
by \((u_0,\nu,f)\).  If \(T_*<\infty\), then it admits a classical extension
beyond \(T_*\), contradicting maximality.
\end{corollary}

\begin{proof}
Theorem~\ref{thm:datum-intersection} supplies the compatible
whole-terminal mixed jet, and Theorem~\ref{thm:datum-rectangles} projects that
record through the independently derived complete-equation rectangle
identities.  The chain from
Theorem~\ref{thm:compatible-inverse-intersection} through
Theorem~\ref{thm:ier-same-history-restart} produces a pressure-normalized classical
extension to \([0,T_*+\delta_0/2]\), where \(\delta_0>0\).  This contradicts
maximality in Definition~\ref{def:formal-maximal-history}.
\end{proof}

\begin{samepage}
\begin{theorem}[Global smoothness and uniqueness for prescribed smooth data]
\label{thm:ier-global-smoothness}
For every \(\nu>0\), every \(u_0\in\HsSigma(\R^3)\), and every prescribed
\(f\in C^\infty([0,\infty);\mathcal S(\R^3;\R^3))\), there exists a
pressure-normalized pair
\begin{equation}
 (u,p)\in C^\infty([0,\infty)\times\R^3;\R^3\times\R)
 \label{eq:ier-global-class}
\end{equation}
satisfying
\begin{align}
 \partial_tu+(u\cdot\nabla)u+\nabla p&=\nu\Delta u+f,
 \label{eq:ier-global-momentum}\\
 \nabla\cdot u&=0,
 \qquad u(0)=u_0,
 \label{eq:ier-global-divergence-datum}\\
 p(t)&=R_iR_j(u_i(t)u_j(t))-(-\Delta)^{-1}\nabla\cdot f(t),
 \qquad \lim_{|x|\to\infty}p(t,x)=0,
 \label{eq:ier-global-pressure}\\
 \frac12\norm{u(t)}_{L^2}^2
 +\nu\int_0^t\norm{\nabla u(s)}_{L^2}^2\dd s
 &=\frac12\norm{u_0}_{L^2}^2
 +\int_0^t\ip{f(s)}{u(s)}_{L^2}\dd s,
 \qquad t\ge0.
 \label{eq:ier-global-energy}
\end{align}
In particular, for every finite \(t\),
\begin{equation}
 \norm{u(t)}_{L^2(\R^3)}
 \le\norm{u_0}_{L^2}+\int_0^t\norm{\Pp f(s)}_{L^2}\dd s<\infty.
 \label{eq:ier-global-energy-bound}
\end{equation}
For every finite \(T>0\), its restriction is the unique member of
\(\mathscr M_{\nu,f}^3(0,T;u_0)\).
\end{theorem}
\end{samepage}

\begin{proof}
 Let
 \[
  \mathscr T:=\left\{\tau>0:
  \begin{array}{c}
  \text{there exists a pressure-complete classical history}\\
  (\tau,u^\tau,p^\tau)\text{ generated by }(u_0,\nu,f)
  \end{array}
  \right\},
 \]
 and put
\[
 T_*:=\sup\mathscr T.
\]
 Theorem~\ref{thm:local-restart}, applied to \(u_0\in\HsSigma\), supplies a
 shifted mild solution whose persistence statement gives every finite
 velocity and pressure row.  Sobolev embedding and the displayed recurrences
 make this pair a pressure-complete classical history, so \(\mathscr T\) is
 nonempty.
 Lemma~\ref{lem:shifted-mild-restriction} admits each strict classical
restriction to the shifted mild class, whose at-most-one property identifies
any two such solutions on their common interval; their union defines a unique
 solution on \([0,T_*)\).  Since \(u_0\in\HsSigma\), the persistence conclusion
of Theorem~\ref{thm:local-restart}, applied on overlapping local intervals,
places every strict restriction in the class of
Definition~\ref{def:formal-maximal-history}.  An extension beyond \(T_*\)
would restrict to an interval \([0,\tau]\) with \(\tau>T_*\), contrary to
the definition of the supremum.  Thus \((T_*,u,p)\) is maximal.  Corollary
\ref{cor:ier-maximal-time-contradiction} gives \(T_*=\infty\).

Fix \(T<\infty\).  Persistence in Theorem~\ref{thm:local-restart}, iterated
along the unique history, retains every finite spatial Sobolev order on
\([0,T]\).  Proposition~\ref{prop:formal-mixed-recurrences} retains every
finite temporal row.  Lemma~\ref{lem:sobolev-embedding} gives
\eqref{eq:ier-global-class}; arbitrariness of \(T\) gives the global class.
The pressure belongs to \(H^s(\R^3)\) for every \(s\), and the Fourier
inversion proof of Lemma~\ref{lem:sobolev-embedding} places its continuous
representative in \(C_0(\R^3)\), which proves the limit in
\eqref{eq:ier-global-pressure}.  Lemma
\ref{lem:shifted-mild-restriction} places \((u,p)|_{[0,T]}\) in
\(\mathscr M_{\nu,f}^3(0,T;u_0)\).  If
\((\widetilde u,\widetilde p)\in\mathscr M_{\nu,f}^3(0,T;u_0)\), the at-most-one
assertion in Theorem~\ref{thm:local-restart} gives
\((\widetilde u,\widetilde p)=(u,p)\) on \([0,T]\).  This proves the stated
uniqueness.

Proposition~\ref{prop:kinetic-energy} holds on every finite interval.  Since
\(T\) was arbitrary, it gives \eqref{eq:ier-global-energy} for every
\(t\ge0\); its regularized square-root estimate gives
\eqref{eq:ier-global-energy-bound}.
\end{proof}
